\section{Policy}
\begin{frame}{}
    \LARGE Reinforcement Learning: \textbf{Policy}
\end{frame}

\begin{frame}{Policy}
\begin{itemize}
    \item Policy $\pi$ determines how the agent chooses actions
    \item $\pi : \mathcal{S} \rightarrow \mathcal{A}$, mapping from states to actions
    \item Deterministic Policy: $$\pi(s) = a$$
    \item Stochastic Policy: $$\pi(a|s) = Pr(a_t = a | s_t = s)$$
\end{itemize}
    
\end{frame}

\begin{frame}{Parameterized Policies}
\begin{itemize}
    \item In tabular settings, we can store $\pi(a|s)$ for every state explicitly
    \item In large or continuous state spaces this is infeasible --- the \textbf{curse of dimensionality}: the number of states grows explosively (e.g., raw-pixel or continuous-control problems), so a table is impossible and we must approximate
    \item Instead, represent the policy as a parameterized function: $\pi_\theta(a|s)$
    \item $\theta$ are learnable parameters (e.g., weights of a neural network)
    \item The network takes state $s$ as input and outputs a probability distribution over actions
    \item Learning RL then becomes: find $\theta^*$ that maximizes expected return
    \item Two broad approaches:
    \begin{itemize}
        \item Learn a value function and derive the policy from it (Days 2, 5, 6)
        \item Directly optimize $\theta$ to maximize return --- Policy Gradient methods (Day 3 onwards)
    \end{itemize}
\end{itemize}
\end{frame}